% Part: first-order-logic 
% Chapter: axiomatic-deduction 
% Section: proving-things-quant

\documentclass[../../../include/open-logic-section]{subfiles}

\begin{document}

\olfileid{fol}{axd}{prq}

\olsection{\usetoken{P}{derivation} met kwantoren}

\begin{ex}
Laten we !!a{derivation} geven van $(\lforall[x][!A(x)] \land
\lforall[y][!B(y)]) \lif \lforall[x][(!A(x) \land !B(x))]$.

Merk eerst op dat
\begin{align*}
  (\lforall[x][!A(x)] \land \lforall[y][!B(y)]) & \lif \lforall[x][!A(x)]\\
  \intertext{een instantie is van \olref[prp]{ax:land1}, en dat}
  \lforall[x][!A(x)] & \lif !A(a) \\
  \intertext{een instantie is van \olref[qua]{ax:q1}. Uit \olref[pro]{prop:chain} weten we dus dat}
  (\lforall[x][!A(x)] \land \lforall[y][!B(y)]) & \lif !A(a) \\
  \intertext{!!{derivable} is. Evenzo geldt dat}
  (\lforall[x][!A(x)] \land \lforall[y][!B(y)]) & \lif \lforall[y][!B(y)] \qquad\text{and}\\
    \lforall[y][!B(y)] & \lif !B(a)\\
    \intertext{respectievelijk instanties zijn van \olref[prp]{ax:land2} en \olref[qua]{ax:q1}, zodat}
    (\lforall[x][!A(x)] \land \lforall[y][!B(y)]) & \lif !B(a)\\
    \intertext{volgens \olref[pro]{prop:chain} afleidbaar is. Met een passende instantie van \olref[prp]{ax:land3} en twee toepassingen van~\MP zien we dat}
    (\lforall[x][!A(x)] \land \lforall[y][!B(y)]) & \lif (!A(a) \land !B(a))\\
    \intertext{afleidbaar is. We kunnen nu \QR{} toepassen en verkrijgen}
(\lforall[x][!A(x)] \land \lforall[y][!B(y)]) & \lif \lforall[x][(!A(x) \land !B(x))].
\end{align*}
\end{ex}

\end{document}
